% chapters/10/section-04.tex
\section{Merekonstruksi dan Memvisualisasikan Jalur}
\label{sec:rekonstruksi}

Pencarian yang berhasil menemukan tujuan baru menyelesaikan separuh pekerjaan.
Bagian ini menyelesaikan separuhnya lagi, sesuai indikator \ind{c}.

\subsection{Menyimpan Asal Setiap Petak}

Selama pencarian, program menyimpan dari petak mana setiap petak ditemukan.
Simpanan itu disebut tabel asal (\istilah{came\_from}), dan bentuknya kamus
yang memetakan sebuah petak ke petak sebelumnya.

\begin{kodekeluaran}
asal[tetangga_petak] = sekarang
\end{kodekeluaran}

Satu baris itu cukup untuk merekonstruksi jalur, sebab setiap petak hanya
memiliki satu asal: yaitu petak yang dilalui untuk mencapainya dengan biaya
termurah. Apabila jalur yang lebih murah ditemukan kemudian, nilai asalnya
ditimpa. Dengan begitu, tabel itu selalu menyatakan jalur termurah yang
diketahui saat itu.

\subsection{Menelusuri Kembali dari Tujuan}

Jalur diperoleh dengan menelusuri tabel asal dari titik tujuan mundur ke titik
awal, lalu membalik urutannya.

\begin{kodekeluaran}
jalur = []
if tujuan in asal:
    petak = tujuan
    while petak is not None:
        jalur.append(petak)
        petak = asal[petak]
    jalur.reverse()
\end{kodekeluaran}

Penelusuran itu berhenti ketika asalnya bernilai \code{None}, yaitu tanda
bahwa petak itu adalah titik awal. Jalur yang dihasilkan
berupa daftar petak dari titik awal ke titik tujuan, sehingga dapat langsung
dipakai NPC sebagai rangkaian tujuan yang harus dilaluinya satu per satu.

\begin{peringatanbox}
Pemeriksaan \code{if tujuan in asal} perlu dilakukan lebih dahulu. Apabila
tujuan tidak dapat dicapai---misalnya karena terkurung dinding---tabel asal
tidak memuat tujuan, dan penelusuran tanpa pemeriksaan itu akan berputar tanpa
henti atau menimbulkan kesalahan. Keadaan tidak dapat dicapai itu wajar terjadi
pada permainan, sehingga harus ditangani, bukan diabaikan.
\end{peringatanbox}

\subsection{Memeriksa Jalur sebelum Dipakai}

Jalur yang dihasilkan sebaiknya diperiksa lebih dahulu, walaupun pencariannya
sudah benar. Pemeriksaan itu murah, dan menangkap kesalahan yang berasal dari
peta atau dari rekonstruksi.

\begin{kodekeluaran}
def periksa_jalur(peta, awal, tujuan, jalur):
    temuan = []
    if not jalur:
        temuan.append("jalur kosong: tujuan tidak dapat dicapai")
    if jalur and jalur[0] != awal:
        temuan.append("jalur tidak dimulai dari titik awal")
    if jalur and jalur[-1] != tujuan:
        temuan.append("jalur tidak berakhir di titik tujuan")
    for petak in jalur:
        if petak in dinding:
            temuan.append("jalur melewati dinding")
    for sebelum, sesudah in zip(jalur, jalur[1:]):
        if jarak_manhattan(sebelum, sesudah) != 1:
            temuan.append("jalur melompat")
    return temuan
\end{kodekeluaran}

Pemeriksaan itu menjawab empat pertanyaan: apakah jalurnya ada, apakah mulainya
benar, apakah akhirnya benar, apakah petaknya bersambung dan tidak menembus
dinding. Pada program bab ini, hasil pemeriksaan ditampilkan pada terminal
setiap kali pencarian dijalankan.

\subsection{Memvisualisasikan Pencarian}

Visualisasi pencarian berguna bukan hanya untuk laporan, melainkan juga untuk
memahami cara kerja algoritmanya. Program bab ini menggambar empat hal.

\begin{itemize}
  \item \textbf{Petak yang telah diperiksa}, berwarna biru gelap. Kumpulan
        petak itu memperlihatkan ke arah mana pencarian menyebar.
  \item \textbf{Petak yang menunggu diperiksa}, berwarna biru terang. Petak itu
        disebut tepi pencarian, dan bentuknya memperlihatkan di mana pencarian
        akan berlanjut.
  \item \textbf{Petak yang sedang diperiksa}, ditandai bingkai kuning. Tanda itu
        memperlihatkan urutan pemeriksaan langkah demi langkah.
  \item \textbf{Jalur yang ditemukan}, berupa bulatan kuning. Jalur itu baru
        digambar setelah pencarian selesai.
\end{itemize}

Pencarian diperlihatkan langkah demi langkah, beberapa petak setiap bingkai,
sehingga penyebarannya dapat diamati. Jumlah petak yang diperiksa setiap
bingkai dapat diubah melalui pilihan \code{--langkah}; nilai besar membuat
pencarian selesai lebih cepat tetapi penyebarannya kurang terlihat.

\begin{tipsbox}
Visualisasi semacam ini juga berguna untuk memeriksa heuristik. Apabila
pencarian menyebar ke arah yang salah---misalnya ke arah yang menjauhi
tujuan---kemungkinan besar heuristiknya keliru. Pada peta yang mengizinkan
delapan arah, penyebaran yang berbentuk bintang menandakan bahwa heuristiknya
melebihkan jarak, dan itulah tanda bahwa jalur yang dihasilkan mungkin bukan
yang terpendek.
\end{tipsbox}
